\chapter{INTRODUÇÃO}

A transformação digital deixou de ser exclusividade de grandes corporações e passou a ser condição de competitividade também para empresas de médio porte, que precisam integrar processos, reduzir retrabalho e responder rapidamente aos clientes \cite{pressman2021}. Em operações logísticas, essa necessidade aparece de forma concreta: cada caixa recebida, guardada, carregada e entregue precisa ser localizada a qualquer momento, e cada veículo precisa ter sua situação conhecida.

Este trabalho tem como cenário uma distribuidora fictícia que recebe mercadorias de fornecedores, armazena as caixas e as entrega aos clientes com frota própria. A empresa já utiliza um ERP para cadastro de clientes, notas fiscais e financeiro, mas o controle do dia a dia --- conferência da entrada, montagem dos carregamentos, liberação de veículos na portaria, registro das entregas e tratamento das devoluções --- é feito em papel, planilhas e mensagens de celular. O resultado é a falta de rastreabilidade: caixas que somem, veículos que saem com carga incompleta e pedidos cuja situação ninguém sabe informar.

O objetivo geral é propor e desenvolver uma solução tecnológica integrada que dê rastreabilidade às caixas e aos veículos, do momento em que entram na empresa até o momento em que são entregues ou retornam. Os objetivos específicos são: caracterizar a organização e seus problemas; planejar a solução e seu modelo de negócio; incorporar ações de inclusão e acessibilidade; desenvolver uma aplicação web, uma API REST e um aplicativo móvel integrados; definir a arquitetura de software e o banco de dados; propor a infraestrutura em nuvem com práticas de DevOps; e organizar o projeto com métodos ágeis.

A metodologia adotada foi a pesquisa aplicada, na forma de um estudo de caso sobre uma empresa fictícia, conforme permitido pelo manual do PIM IV. O processo atual (\textit{as-is}) foi descrito e analisado para levantar requisitos e regras de negócio; em seguida, a solução foi construída de forma incremental, em ciclos curtos, com testes automatizados executados ao final de cada incremento \cite{sommerville2018}. As decisões técnicas foram fundamentadas na bibliografia das disciplinas do semestre: Empreendedorismo em TI; Relações Étnico-Raciais e Ensino de Culturas Afro-Indígenas; Desenvolvimento Web com .NET; Desenvolvimento Mobile; Programação Aplicada em .NET; Programação de Banco de Dados; Cloud Computing e DevOps; e Gerenciamento de Projetos Ágil.

Para validação e reprodutibilidade, o código-fonte completo (web, API, aplicativo, testes e este relatório) está publicado em \url{https://github.com/Joao-Pedro-ADS/PIM-IV-Logistica}. O relatório segue a ordem das etapas do manual. O Capítulo 2 caracteriza a organização e o Capítulo 3 apresenta o planejamento da solução, o modelo de negócio e os requisitos. O Capítulo 4 trata de responsabilidade social e inclusão. Os Capítulos 5 e 6 descrevem as soluções web e móvel; os Capítulos 7 e 8, a arquitetura de software e o banco de dados. O Capítulo 9 apresenta a proposta de nuvem e DevOps, o Capítulo 10 o gerenciamento ágil e o Capítulo 11 os resultados dos testes, seguidos da conclusão.
